\section{The four graph codes}

Throughout, automorphisms of a code mean permutations of its coordinates,
not arbitrary monomial transformations. Unless a prime is displayed, the
field is $F=\mathbb F_3$. The constant vector is denoted by $j$, with its
length determined by context. A direct sum of submodules below is an
internal sum in the same coordinate space, rather than a concatenation of
shorter codes. For a symmetric adjacency matrix $A_n$ of order $n$, put
\[
 C_n=\im A_n,\qquad
 D_n=\langle (A_n)_i-(A_n)_1:1\leq i\leq n\rangle.
\]
For length $120$ we also write $S=D_{120}$, to agree with the source.
The code of the complementary graph is denoted by $\overline C_n$.
This notation distinguishes a graph complement from an orthogonal dual.

The main results of Leemans and Rodrigues~\cite{LR} concern four specified
strongly regular graphs. We use their graph constructions and the
identifications with the relevant McLaughlin subgraphs, as attributed in
that paper to Brouwer--van Eijl and the preceding graph literature.
The actual graphs are fixed here, so assertions about a particular code
are not inferred merely from four strongly regular parameters.

For $n=105$, use the incident point--line flags of $\operatorname{PG}(2,4)$.
Join distinct flags $(p,\ell),(q,m)$ if $p=q$, or $\ell=m$, or neither
cross incidence $p\in m$, $q\in\ell$ holds. The finite field is
$\mathbb F_4=\mathbb F_2[z]/(z^2+z+1)$. All $105$ vertices and their
adjacencies are constructed directly from this field. The printed
adjacency sentence on page~242 has a defect: a conjunctive reading gives
valency eight, while joining all its clauses disjunctively gives forty.
Our identification with the intended graph therefore retains Coolsaet's
uniqueness theorem for the complementary parameters $(105,32,4,12)$
\cite[Theorem~3]{Coolsaet}; the verified complementary parameters alone
would not identify a graph without that input.

For the other three lengths, let $\mathcal D$ be the specified
$S(4,7,23)$ with its $253$ heptads. One complete construction takes the
weight-seven words in the binary span of the twelve shifts of
\[
 1+x+x^5+x^6+x^7+x^9+x^{11}.
\]
The supplied constructor verifies all $4096$ words, their weight
multiplicities, and the unique containment of every four-subset in a
heptad. It also verifies all block intersection sizes. Delete a fixed set
of $t$ points, where $t=0,1,2$, and retain the heptads avoiding it. There
are respectively $253,176,120$ such blocks. Join distinct retained blocks
when their intersection has size three. This is exactly the complement
of the source's intersection-one graph. With $H_t$ the point--block
incidence matrix on the remaining $23-t$ points, the common paper
\cite[Theorem~2.2]{U} gives
\[
\begin{array}{c|c|c|c}
 t&n&H_tH_t^{\mathsf T}\pmod3&C_n\\\hline
 0&253&2I_{23}&\im H_0^{\mathsf T}\\
 1&176&I_{22}+J_{22}&H_1^{\mathsf T}(e^\perp)\\
 2&120&I_{21}&\im H_2^{\mathsf T}.
\end{array}                                                    \tag{1}
\]
Here $e$ is the constant point vector. Each of these three codes is LCD;
$\overline C_n=C_n^\perp$ for $n=120,176$. Also
$C_{253}=Fj\mathbin{\perp\!\oplus}D_{253}$, with
$D_{253}=H_0^{\mathsf T}(e^\perp)$ of dimension $22$ and LCD.
The hypotheses of the common theorem are discharged by the complete
Steiner and intersection checks, not by assuming the ternary ranks.

The verified parameters and integral spectra are
\[
\begin{array}{c|c|c}
n&(n,k,\lambda,\mu)&\text{spectrum of }A_n\\\hline
105&(105,72,51,45)&72^1,9^{20},(-3)^{84}\\
120&(120,77,52,44)&77^1,11^{20},(-3)^{99}\\
176&(176,105,68,54)&105^1,17^{21},(-3)^{154}\\
253&(253,140,87,65)&140^1,25^{22},(-3)^{230}.
\end{array}                                                     \tag{2}
\]
The finite certificate checks every vertex degree and every pair's
common-neighbour count. The spectra then follow from
$A^2=(\lambda-\mu)A+(k-\mu)I+\mu J$, the constant eigenvector,
and the trace and dimension identities. Connectivity follows from
$\mu>0$, so the constant eigenvalue has multiplicity one.

\subsection{The retained orbital construction}
The source's Result~1 attributes a general construction to Key and Moori,
including their correction. Its argument is short enough to make the
exact scope explicit. Let $G$ act primitively on $\Omega$, let $\alpha\in
\Omega$, and let $\Delta\ne\{\alpha\}$ be a nonempty orbit of $G_\alpha$.
The set stabilizer $G_\Delta$ contains $G_\alpha$. Primitivity makes the
latter maximal; since $\Delta$ is a proper nonempty subset,
$G_\Delta\ne G$. Hence $G_\Delta=G_\alpha$, and the translates of $\Delta$
are $|\Omega|$ distinct blocks. Transitivity and counting incidences give
constant block size and replication $|\Delta|$, producing the symmetric
$1$-design. If the orbital is self-paired, these blocks are precisely the
neighbour sets of the associated undirected graph. The nonempty graph's
components form a $G$-invariant partition, so it is connected. Its
symmetric incidence matrix identifies the design with its dual, and both
point and block actions are primitive. This preserves the background
result without using it in place of any of the explicit constructions above.
